(1) 圆柱对称性使电场仅有环向分量 \(E_\phi(r,t)\)。由 Maxwell–Faraday 积分定律
\[
2\pi rE_\phi=-\frac{d\Phi_B(r)}{dt}.
\]
当 r<a，所围磁通为 \(\pi r^2B(t)\)，所以 \(E_\phi=-r\dot B/2\)；当 r>a，磁通为 \(\pi a^2B(t)\)，所以
\[
E_\phi(r,t)=-\frac{a^2}{2r}\dot B(t).
\]
两个公式在 r=a 连续，不要求管外磁感应强度非零。

(2) 对 r<a，环量 \(2\pi rE_\phi=-\pi r^2\dot B\)；对 r>a，环量 \(2\pi rE_\phi=-\pi a^2\dot B\)。正好等于各圆盘面的穿透磁通变化率的负值。管外 E 对 r 的依赖为 1/r，故不同的外部圆周都环绕同一磁通，其电场环量彼此相同。

(3) 电荷沿圆周受切向电力 \(qE\varphi \)。当在足够短于 B 显著变化的时间里绕行一周，可将 \(\dot B\) 当作常数，电场做功
\[
W_E=q\oint\boldsymbol E\cdot d\boldsymbol\ell
=-q\pi a^2\dot B.
\]
若时间不可忽略则必须沿真实时空轨道积分 \(q\int E_\phi[r,t(\phi)]r\,d\phi\)，不能把不同时间的场抽成某一固定瞬间的环量。外力保持匀速时其功补偿电场功（导轨无摩擦），加速的能量来源于外部改变磁通的装置。

(4) 外部区域局域旋度 \(\nabla\times\boldsymbol E=0\)（磁场变化局限在管内），但绕管的一条闭曲线不能在不穿过磁场区域的情况下缩成一点；其非零环量体现整体拓扑。静态 \(\dot B=0\) 时所有感应电场均为零，不能将外部静磁通本身当成永久感应电动势。
